\section{Fazit und Ausblick}
\label{sec:conclusion}

Schreibaktivität lässt sich aus den IMU-Daten einer Apple Watch am
Schreibhandgelenk personenunabhängig erkennen. Mit 32 Personen, einem
instrumentierten Stift als Ground Truth und einer Evaluation ohne Person im
eigenen Training erreicht ein Faltungsnetz mit rekurrentem Kopf auf
5-s-Fenstern 0,921 Accuracy und 0,975 ROC-AUC (drei Seeds, ohne gepaarten
Test). Ein Random Forest auf 88 Merkmalen erreicht je Sekunde 0,865, mit
kausalem HMM 0,898 und rückblickend 0,917; das HMM läuft ohne Neutraining im
Live-System.

Drei Erkenntnisse gehen über die Zahlen hinaus. Erstens liegt die Grenze des
Merkmalsmodells eher in der Repräsentation als im Signal: Zwölf zusätzliche
Personen haben den Random Forest nicht bewegt, die Netze schon. Zweitens ist
Zeitkontext der wirksamste Hebel und lässt sich nur einmal einsammeln, mit
zwei Prozentpunkten Aufschlag für alles, was nicht live entschieden werden
muss. Drittens sitzen die verbleibenden Fehler bei bestimmten Personen in
bestimmten Aufgaben, vor allem beim Tippen, und dort ist Trainingsvielfalt
der Hebel, nicht Merkmalsdesign.

Der nächste Schritt liegt deshalb in den Daten: Personen aufnehmen, die anders
tippen und schreiben als die bisherige Kohorte, und einen Teil der Aufnahmen
gegen Video annotieren, damit die Labeldefinition selbst überprüfbar wird.
Beides verspricht mehr als weitere Rechenzeit.
